% =============================================================================
\chapter{Giải thuật}\label{chap:giaithuat}
% =============================================================================

Ký hiệu dùng trong chương: $n$ -- số mặt hàng; $m$ -- bậc B-Tree; $t = \lceil m/2\rceil$;
$h$ -- chiều cao cây; $L$ -- tổng số lô; $P$ -- số phiếu; $r$ -- số phần tử kết quả.

% -----------------------------------------------------------------------------
\section{B-Tree bậc $m$}\label{sec:btree}
% -----------------------------------------------------------------------------

\subsection{Định nghĩa và cách lưu nút}

B-Tree bậc $m$~\cite{bayer,clrs} là cây tìm kiếm nhiều nhánh thỏa mãn:
(1) mỗi nút có tối đa $m-1$ khóa và $m$ con; (2) mỗi nút khác gốc có ít nhất $t-1 = \lceil m/2\rceil-1$ khóa;
gốc có ít nhất 1 khóa; (3) các khóa trong nút tăng dần, và mọi khóa trong cây con thứ $i$ nằm giữa khóa
$i-1$ và khóa $i$ của nút cha; (4) mọi lá ở cùng một mức.

Mỗi nút cấp phát mảng \code{keys} có $m$ ô và mảng \code{child} có $m+1$ ô -- \emph{dư một ô} so với định nghĩa.
Ô dư cho phép thuật toán chèn đặt khóa vào nút trước, để nút tạm thời có $m$ khóa (tràn), rồi mới tách. Cách
``chèn rồi tách ngược lên'' này đúng với mọi $m \ge 3$, kể cả $m$ lẻ như 3 và 5; còn cách tách phòng ngừa
khi đi xuống trong~\cite{clrs} chỉ áp dụng cho bậc chẵn $m = 2t$.

\subsection{Tìm kiếm}

Trong mỗi nút, vị trí khóa được tìm bằng tìm nhị phân (\textsc{TimViTri} trả về chỉ số $i$ nhỏ nhất có
$\mathit{keys}[i] \ge k$). Mỗi nút đi qua được cộng vào bộ đếm \code{soNutTruyXuat} -- phục vụ phần nâng cao.

\begin{algorithm}[htbp]
\caption{Tìm mặt hàng theo mã vạch (\code{btSearch})}\label{alg:tim}
\begin{algorithmic}[1]
\Function{Tim}{$T, k$}
  \State $x \gets T.\mathit{root}$
  \While{$x \neq \textsc{nil}$}
    \State $T.\mathit{soNutTruyXuat} \gets T.\mathit{soNutTruyXuat} + 1$
    \State $i \gets \Call{TimViTri}{x, k}$ \Comment{tìm nhị phân, $O(\log m)$}
    \If{$i < x.n$ \textbf{và} $x.\mathit{keys}[i].\mathit{maVach} = k$} \State \Return $x.\mathit{keys}[i]$ \EndIf
    \If{$x.\mathit{leaf}$} \State \Return \textsc{nil} \EndIf
    \State $x \gets x.\mathit{child}[i]$
  \EndWhile
  \State \Return \textsc{nil}
\EndFunction
\end{algorithmic}
\end{algorithm}

\textbf{Độ phức tạp:} đi qua tối đa $h$ nút, mỗi nút $O(\log m)$ phép so sánh chuỗi:
$O(h\log m) = O(\log n)$ phép so sánh và $O(h) = O(\log_t n)$ nút.

\subsection{Chèn và tách nút}

\begin{algorithm}[htbp]
\caption{Chèn mặt hàng (\code{btInsert}) và tách nút}\label{alg:chen}
\begin{algorithmic}[1]
\Function{Chen}{$T, \mathit{mh}$}
  \If{$T.\mathit{root} = \textsc{nil}$} \State tạo lá gốc chứa $\mathit{mh}$; \Return \textbf{đúng} \EndIf
  \If{\textbf{không} \Call{ChenDeQuy}{$T, T.\mathit{root}, \mathit{mh}$}} \State \Return \textbf{sai} \Comment{trùng mã} \EndIf
  \If{$T.\mathit{root}.n = m$} \Comment{gốc tràn: cây cao thêm 1 mức}
    \State $s \gets$ nút trong mới; $s.\mathit{child}[0] \gets T.\mathit{root}$; $T.\mathit{root} \gets s$
    \State \Call{TachCon}{$T, s, 0$}
  \EndIf
  \State $T.\mathit{size} \gets T.\mathit{size} + 1$; \Return \textbf{đúng}
\EndFunction
\Statex
\Function{ChenDeQuy}{$T, x, \mathit{mh}$}
  \State $i \gets \Call{TimViTri}{x, \mathit{mh}.\mathit{maVach}}$
  \If{$i < x.n$ \textbf{và} $x.\mathit{keys}[i].\mathit{maVach} = \mathit{mh}.\mathit{maVach}$} \State \Return \textbf{sai} \EndIf
  \If{$x.\mathit{leaf}$}
    \State dời $x.\mathit{keys}[i..x.n-1]$ sang phải 1 ô; $x.\mathit{keys}[i] \gets \mathit{mh}$; $x.n \gets x.n+1$
    \State \Return \textbf{đúng} \Comment{$x$ có thể đang có $m$ khóa}
  \EndIf
  \If{\textbf{không} \Call{ChenDeQuy}{$T, x.\mathit{child}[i], \mathit{mh}$}} \State \Return \textbf{sai} \EndIf
  \If{$x.\mathit{child}[i].n = m$} \State \Call{TachCon}{$T, x, i$} \EndIf
  \State \Return \textbf{đúng}
\EndFunction
\Statex
\Procedure{TachCon}{$T, x, i$} \Comment{$y = x.\mathit{child}[i]$ đang có $m$ khóa}
  \State $y \gets x.\mathit{child}[i]$; $g \gets \lfloor m/2 \rfloor$
  \State $z \gets$ nút mới, $z.\mathit{leaf} \gets y.\mathit{leaf}$, $z.n \gets m-1-g$
  \State $z.\mathit{keys} \gets y.\mathit{keys}[g+1..m-1]$; \textbf{nếu} $y$ không là lá: $z.\mathit{child} \gets y.\mathit{child}[g+1..m]$
  \State $y.n \gets g$
  \State chèn $y.\mathit{keys}[g]$ vào $x.\mathit{keys}[i]$ và $z$ vào $x.\mathit{child}[i+1]$ (dời phải)
  \If{$T.\mathit{inKhiTachGop}$} \State in thông báo tách và (sau thao tác) in cả cây \EndIf
\EndProcedure
\end{algorithmic}
\end{algorithm}

\paragraph{Tính đúng của phép tách.} Nút tràn có $m$ khóa; khóa giữa $\mathit{keys}[g]$ lên cha, nút trái giữ
$g = \lfloor m/2\rfloor$ khóa, nút phải giữ $m-1-\lfloor m/2\rfloor = \lceil m/2\rceil - 1 = t-1$ khóa. Cả hai
đều đạt số khóa tối thiểu. Cha nhận thêm một khóa và có thể tràn tiếp; khi đó lời gọi đệ quy ở mức trên tách
cha, lan dần tới gốc.

\paragraph{Ví dụ chạy tay ($m = 3$).} Mỗi nút có 1--2 khóa. Chèn lần lượt 50, 20, 80, 10, 30, 60, 90, 40, 70, 25
(viết gọn mã vạch \code{00000050} thành 50). Hình~\ref{fig:chen} là cây sau các bước có tách nút; nút tô màu là
nút vừa được tạo hoặc thay đổi. Các bước này trùng khớp với kết quả chương trình in ra ở chế độ
\code{./test\_btree demo} (Mục~\ref{sec:manhinh}).

\forestset{bt/.style={for tree={draw,rounded corners=1pt,font=\ttfamily\footnotesize,inner sep=1.5pt,
  s sep=1.5mm,l sep=4mm,edge={thick,black!60}}},
  moi/.style={fill=orange!25}}

\begin{figure}[htbp]
\centering
\begin{subfigure}[b]{0.22\linewidth}\centering
\begin{forest} bt [50,moi [20,moi] [80,moi]] \end{forest}
\caption{Chèn 50, 20, 80: lá \code{[20|50|80]} tràn, tách, 50 lên gốc mới}
\end{subfigure}\hfill
\begin{subfigure}[b]{0.25\linewidth}\centering
\begin{forest} bt [{20 | 50},moi [10,moi] [30,moi] [80]] \end{forest}
\caption{Chèn 10, 30: lá \code{[10|20|30]} tràn, 20 lên cha}
\end{subfigure}\hfill
\begin{subfigure}[b]{0.24\linewidth}\centering
\begin{forest} bt [50,moi [20,moi [10] [30]] [80,moi [60,moi] [90,moi]]] \end{forest}
\caption{Chèn 60, 90: tách lá, 80 lên làm gốc tràn, tách gốc: cây cao 3}
\end{subfigure}\hfill
\begin{subfigure}[b]{0.26\linewidth}\centering
\begin{forest} bt [50 [{20 | 30},moi [10] [25,moi] [40,moi]] [80 [{60 | 70}] [90]]] \end{forest}
\caption{Chèn 40, 70, 25: lá \code{[25|30|40]} tràn, 30 lên cha}
\end{subfigure}
\caption{Chèn và tách nút trên B-Tree bậc 3}
\label{fig:chen}
\end{figure}

\subsection{Xóa: mượn và gộp nút}

Xóa một khóa ở nút trong được quy về xóa ở lá: khóa bị thay bằng \emph{khóa liền trước} (khóa lớn nhất của cây
con trái), rồi xóa khóa liền trước đó khỏi lá. Sau khi xóa ở cây con, nếu nút con còn ít hơn $t-1$ khóa (thiếu),
nút cha sửa lại bằng một trong ba cách theo thứ tự ưu tiên: mượn anh trái, mượn anh phải, gộp với anh em.

\begin{algorithm}[htbp]
\caption{Xóa mặt hàng (\code{btRemove}) với mượn và gộp nút}\label{alg:xoa}
\begin{algorithmic}[1]
\Function{Xoa}{$T, k$}
  \State $\mathit{kq} \gets \Call{XoaDeQuy}{T, T.\mathit{root}, k}$
  \If{$\mathit{kq} \neq \textsc{nil}$ \textbf{và} $T.\mathit{root}.n = 0$} \Comment{gốc rỗng: cây thấp đi 1 mức}
    \State $T.\mathit{root} \gets$ (\textbf{nếu} gốc là lá \textbf{thì} \textsc{nil} \textbf{ngược lại} $T.\mathit{root}.\mathit{child}[0]$); giải phóng gốc cũ
  \EndIf
  \State \Return $\mathit{kq}$ \Comment{con trỏ mặt hàng; \code{Inventory} chịu trách nhiệm \code{delete}}
\EndFunction
\Statex
\Function{XoaDeQuy}{$T, x, k$}
  \State $i \gets \Call{TimViTri}{x, k}$
  \If{$i < x.n$ \textbf{và} $x.\mathit{keys}[i].\mathit{maVach} = k$}
    \State $\mathit{kq} \gets x.\mathit{keys}[i]$
    \If{$x.\mathit{leaf}$} \State dời trái $x.\mathit{keys}[i+1..]$; $x.n \gets x.n - 1$; \Return $\mathit{kq}$ \EndIf
    \State $x.\mathit{keys}[i] \gets \Call{XoaMax}{T, x.\mathit{child}[i]}$ \Comment{thay bằng khóa liền trước}
  \ElsIf{$x.\mathit{leaf}$} \State \Return \textsc{nil} \Comment{không có khóa}
  \Else \State $\mathit{kq} \gets \Call{XoaDeQuy}{T, x.\mathit{child}[i], k}$
  \EndIf
  \If{$\mathit{kq} \neq \textsc{nil}$ \textbf{và} $x.\mathit{child}[i].n < t - 1$} \State \Call{SuaThieu}{$T, x, i$} \EndIf
  \State \Return $\mathit{kq}$
\EndFunction
\Statex
\Procedure{SuaThieu}{$T, x, i$} \Comment{$c = x.\mathit{child}[i]$ có $t-2$ khóa}
  \If{$i > 0$ \textbf{và} $x.\mathit{child}[i-1].n > t-1$} \Comment{mượn anh trái}
    \State $c$ nhận $x.\mathit{keys}[i-1]$ ở đầu; $x.\mathit{keys}[i-1] \gets$ khóa cuối của anh trái
    \State con cuối của anh trái thành con đầu của $c$
  \ElsIf{$i < x.n$ \textbf{và} $x.\mathit{child}[i+1].n > t-1$} \Comment{mượn anh phải (đối xứng)}
    \State $c$ nhận $x.\mathit{keys}[i]$ ở cuối; $x.\mathit{keys}[i] \gets$ khóa đầu của anh phải
  \ElsIf{$i > 0$} \State \Call{Gop}{$x, i-1$} \Else \State \Call{Gop}{$x, i$} \EndIf
\EndProcedure
\Statex
\Procedure{Gop}{$x, j$} \Comment{gộp $\mathit{child}[j]$, $\mathit{keys}[j]$, $\mathit{child}[j+1]$}
  \State $\mathit{trai}.\mathit{keys} \gets \mathit{trai}.\mathit{keys} \cdot x.\mathit{keys}[j] \cdot \mathit{phai}.\mathit{keys}$; nối các con tương ứng
  \State xóa $x.\mathit{keys}[j]$ và $x.\mathit{child}[j+1]$ khỏi $x$ (dời trái); giải phóng nút $\mathit{phai}$
\EndProcedure
\end{algorithmic}
\end{algorithm}

\textsc{XoaMax}$(T,x)$ đi theo con phải nhất xuống lá, gỡ khóa cuối, và trên đường quay lên gọi \textsc{SuaThieu}
cho nút con bị thiếu -- giống \textsc{XoaDeQuy}.

\paragraph{Tính đúng của phép gộp.} Chỉ gộp khi anh em có đúng $t-1$ khóa (không mượn được). Nút sau gộp có
$(t-2) + 1 + (t-1) = 2t-2 = 2\lceil m/2\rceil - 2 \le m-1$ khóa, vừa trong một nút. Cha mất một khóa và có thể
thiếu; lời gọi ở mức trên sẽ sửa tiếp. Nếu gốc mất khóa cuối cùng, con duy nhất của nó trở thành gốc mới.

\paragraph{Ví dụ chạy tay.} Tiếp tục từ cây ở Hình~\ref{fig:chen}d, xóa lần lượt 20, 30, 50, 10, 90
(Hình~\ref{fig:xoa}).

\begin{figure}[htbp]
\centering
\begin{subfigure}[b]{0.25\linewidth}\centering
\begin{forest} bt [50 [30,moi [{10 | 25},moi] [40]] [80 [{60 | 70}] [90]]] \end{forest}
\caption{Xóa 20: thay bằng 10; lá rỗng, anh phải \code{[25]} không cho mượn được, gộp thành \code{[10|25]}}
\end{subfigure}\hfill
\begin{subfigure}[b]{0.23\linewidth}\centering
\begin{forest} bt [50 [25,moi [10,moi] [40]] [80 [{60 | 70}] [90]]] \end{forest}
\caption{Xóa 30: thay bằng khóa liền trước 25; lá \code{[10]} không thiếu}
\end{subfigure}\hfill
\begin{subfigure}[b]{0.24\linewidth}\centering
\begin{forest} bt [{40 | 80},moi [{10 | 25},moi] [{60 | 70}] [90]] \end{forest}
\caption{Xóa 50: thay bằng 40, gộp 2 lần, gốc rỗng, cây thấp đi}
\end{subfigure}\hfill
\begin{subfigure}[b]{0.24\linewidth}\centering
\begin{forest} bt [{40 | 70},moi [25] [60,moi] [80,moi]] \end{forest}
\caption{Xóa 10, rồi 90: lá rỗng, mượn anh trái: 70 lên, 80 xuống}
\end{subfigure}
\caption{Xóa trên B-Tree bậc 3: gộp nút, cây thấp đi và mượn nút}
\label{fig:xoa}
\end{figure}

\subsection{Duyệt, đếm và in cây}

\begin{itemize}
  \item \code{btTraverse}: duyệt trung thứ tự (con $0$, khóa $0$, con $1$, khóa $1$,~\ldots), gọi hàm
        \code{visit(mh, ctx)} cho từng mặt hàng theo thứ tự mã vạch tăng dần. Các mô-đun khác dùng hàm này để
        ghi tệp, lập báo cáo và giải phóng bộ nhớ mà không cần biết cấu trúc nút. $O(n)$.
  \item \code{btHeight} đi theo con trái nhất: $O(h)$; \code{btSoNut} đếm đệ quy: $O(n)$.
  \item \code{btPrint} in cây theo từng mức (\code{Muc 0}, \code{Muc 1},~\ldots) bằng cách duyệt lại cây cho
        mỗi mức: $O(n\cdot h)$, chỉ dùng cho cây nhỏ khi trình diễn.
  \item \code{btDestroy} giải phóng mọi nút nhưng \emph{không} giải phóng mặt hàng (theo quy ước sở hữu).
\end{itemize}

% -----------------------------------------------------------------------------
\section{Nhập kho và xuất kho FIFO}
% -----------------------------------------------------------------------------

\begin{algorithm}[htbp]
\caption{Nhập kho và xuất kho theo FIFO (\code{nhapKho}, \code{xuatKho})}\label{alg:fifo}
\begin{algorithmic}[1]
\Function{NhapKho}{$\mathit{kho}, \mathit{ma}, \mathit{sl}, \mathit{giaNhap}, \mathit{hsd}, \mathit{ngay}$}
  \If{$\mathit{sl} \le 0$ \textbf{hoặc} ngày không hợp lệ \textbf{hoặc} $\mathit{hsd} < \mathit{ngay}$} \State \Return \code{KQ\_DU\_LIEU\_SAI} \EndIf
  \State $\mathit{mh} \gets \Call{Tim}{\mathit{kho}.\mathit{cay}, \mathit{ma}}$; \textbf{nếu} $\mathit{mh} = \textsc{nil}$ \Return \code{KQ\_KHONG\_TIM\_THAY}
  \State $\mathit{lo} \gets$ lô mới (mã $\mathit{kho}.\mathit{maLoKeTiep}$++, $\mathit{ngay}$, $\mathit{hsd}$, $\mathit{sl}$, $\mathit{giaNhap}$)
  \State \Call{Enqueue}{$\mathit{mh}.\mathit{dsLo}, \mathit{lo}$} \Comment{thêm vào cuối, $O(1)$}
  \State $\mathit{mh}.\mathit{tonKho} \gets \mathit{mh}.\mathit{tonKho} + \mathit{sl}$; thêm phiếu \code{NHAP} vào cuối danh sách phiếu
\EndFunction
\Statex
\Function{XuatKho}{$\mathit{kho}, \mathit{ma}, \mathit{sl}, \mathit{ngay}$}
  \State kiểm tra $\mathit{sl} > 0$, ngày hợp lệ; $\mathit{mh} \gets \Call{Tim}{\mathit{kho}.\mathit{cay}, \mathit{ma}}$
  \If{$\mathit{mh}.\mathit{tonKho} < \mathit{sl}$} \State \Return \code{KQ\_KHONG\_DU\_HANG} \Comment{kiểm tra trước, không xuất dở dang}\EndIf
  \State $\mathit{can} \gets \mathit{sl}$
  \While{$\mathit{can} > 0$}
    \State $\mathit{lo} \gets \Call{Front}{\mathit{mh}.\mathit{dsLo}}$ \Comment{lô cũ nhất}
    \If{$\mathit{lo}.\mathit{soLuong} \le \mathit{can}$}
      \State $\mathit{can} \gets \mathit{can} - \mathit{lo}.\mathit{soLuong}$; \Call{Dequeue}{$\mathit{mh}.\mathit{dsLo}$}; giải phóng $\mathit{lo}$
    \Else \State $\mathit{lo}.\mathit{soLuong} \gets \mathit{lo}.\mathit{soLuong} - \mathit{can}$; $\mathit{can} \gets 0$
    \EndIf
  \EndWhile
  \State $\mathit{mh}.\mathit{tonKho} \gets \mathit{mh}.\mathit{tonKho} - \mathit{sl}$; thêm phiếu \code{XUAT} với đơn giá bán hiện tại
\EndFunction
\end{algorithmic}
\end{algorithm}

\textbf{Độ phức tạp:} nhập kho $O(\log n)$ (tìm) $+\ O(1)$; xuất kho $O(\log n + k)$ với $k$ là số lô bị lấy
hết -- mỗi lô chỉ bị lấy hết một lần trong đời nên tổng chi phí khấu trừ của mọi lần xuất là $O(\log n)$ mỗi lần.
Xuất kho kiểm tra đủ hàng \emph{trước} khi thay đổi dữ liệu, nên thao tác hoặc thành công trọn vẹn hoặc không
làm thay đổi gì.

\paragraph{Ví dụ chạy tay (dữ liệu mẫu).} Mặt hàng \code{8930000000002} được nhập 3 lô 38, 38, 39 đơn vị
(\code{PN00003}--\code{PN00005}); phiếu \code{PX00006} xuất 51 đơn vị. Lô 3 (38) bị lấy hết và gỡ khỏi hàng đợi,
còn thiếu $51-38 = 13$ đơn vị lấy từ lô 4, lô 4 còn 25 (Hình~\ref{fig:fifo}). Kết quả khớp với \code{lo.txt}.

\begin{figure}[htbp]
\centering
\begin{tikzpicture}[
  lo/.style={rectangle split,rectangle split parts=2,draw,font=\scriptsize\ttfamily,align=center,
             minimum width=1.7cm,inner sep=2pt,fill=cTan!60},
  het/.style={lo,fill=black!8,text=black!45,dashed},
  ptr/.style={-{Stealth[length=1.8mm]},thick}]
  \node[font=\small\bfseries,anchor=east] at (-0.6,0) {Trước};
  \node[lo] (a3) at (1,0) {LO3\nodepart{two}38};
  \node[lo] (a4) at (3.2,0) {LO4\nodepart{two}38};
  \node[lo] (a5) at (5.4,0) {LO5\nodepart{two}39};
  \draw[ptr] (a3) -- (a4); \draw[ptr] (a4) -- (a5);
  \node[font=\scriptsize,anchor=west] at (6.4,0) {tồn = 115};
  \node[font=\scriptsize] (f1) at (1,1) {front}; \draw[ptr] (f1) -- (a3);
  \node[font=\scriptsize] (r1) at (5.4,1) {rear}; \draw[ptr] (r1) -- (a5);
  \node[font=\small\bfseries,anchor=east] at (-0.6,-2) {Sau};
  \node[het] (b3) at (1,-2) {LO3\nodepart{two}0};
  \node[lo,fill=orange!30] (b4) at (3.2,-2) {LO4\nodepart{two}25};
  \node[lo] (b5) at (5.4,-2) {LO5\nodepart{two}39};
  \draw[ptr] (b4) -- (b5);
  \node[font=\scriptsize,anchor=west] at (6.4,-2) {tồn = 64};
  \node[font=\scriptsize] (f2) at (3.2,-1) {front}; \draw[ptr] (f2) -- (b4);
  \node[font=\scriptsize,align=center] at (1,-3) {lấy 38, gỡ lô};
  \node[font=\scriptsize,align=center] at (3.2,-3) {lấy 13};
  \draw[-{Stealth},very thick,black!50] (8.3,-0.3) -- node[right,font=\scriptsize,align=left]{xuất 51\\(\code{PX00006})} (8.3,-1.7);
\end{tikzpicture}
\caption{Xuất kho FIFO 51 đơn vị mặt hàng \code{8930000000002}}
\label{fig:fifo}
\end{figure}

% -----------------------------------------------------------------------------
\section{Merge Sort tự cài đặt}
% -----------------------------------------------------------------------------

Merge Sort~\cite{sedgewick} được cài dưới dạng một hàm mẫu (template) \code{mergeSort<T>(a, n, truoc)} dùng chung cho mọi kiểu
sắp xếp; \code{truoc(x, y)} trả về \textbf{đúng} nếu $x$ phải đứng trước $y$. Mảng tạm $n$ phần tử được cấp phát
một lần và dùng lại ở mọi mức đệ quy.

\begin{algorithm}[htbp]
\caption{Merge Sort ổn định (\code{Sort.h})}\label{alg:mergesort}
\begin{algorithmic}[1]
\Procedure{SapDeQuy}{$a, \mathit{tam}, l, r$} \Comment{sắp $a[l..r-1]$}
  \If{$r - l < 2$} \State \Return \EndIf
  \State $g \gets l + \lfloor (r-l)/2 \rfloor$; \Call{SapDeQuy}{$a, \mathit{tam}, l, g$}; \Call{SapDeQuy}{$a, \mathit{tam}, g, r$}
  \State $i \gets l$; $j \gets g$; $k \gets l$
  \While{$i < g$ \textbf{và} $j < r$}
    \If{$\mathit{truoc}(a[j], a[i])$} \State $\mathit{tam}[k{+}{+}] \gets a[j{+}{+}]$ \Comment{chỉ lấy bên phải khi \emph{thật sự} đứng trước}
    \Else \State $\mathit{tam}[k{+}{+}] \gets a[i{+}{+}]$ \Comment{bằng nhau: lấy bên trái $\Rightarrow$ ổn định}
    \EndIf
  \EndWhile
  \State chép phần còn lại của hai nửa vào $\mathit{tam}$; chép $\mathit{tam}[l..r-1]$ về $a$
\EndProcedure
\end{algorithmic}
\end{algorithm}

\textbf{Độ phức tạp:} $T(n) = 2T(n/2) + \Theta(n) \Rightarrow \Theta(n\log n)$ ở mọi trường hợp; bộ nhớ phụ
$\Theta(n)$; độ sâu đệ quy $\lceil\log_2 n\rceil$. Các tiêu chí sắp xếp được dùng:

\begin{table}[htbp]
\centering\small
\caption{Các kiểu sắp xếp dùng trong báo cáo}
\label{tab:sapxep}
\begin{tabular}{@{}lll@{}}
\toprule
\textbf{Hàm} & \textbf{Tiêu chí} & \textbf{Dùng cho} \\
\midrule
\code{sapXepTheoTen} & tên (không phân biệt hoa thường), rồi mã vạch & tìm theo tên, liệt kê 1 nhóm \\
\code{sapXepTheoNhom} & nhóm, rồi tên, rồi mã vạch & liệt kê tất cả, gom nhóm \\
\code{sapXepTheoTonKho} & tồn kho tăng dần & cảnh báo tồn dưới ngưỡng \\
\code{sapXepTheoHSD} & HSD tăng dần & cảnh báo sắp hết hạn \\
\code{sapXepBanChay} & số lượng bán giảm dần, rồi mã vạch & top 10 bán chạy \\
\code{sapXepGiaTriNhom} & giá trị tồn giảm dần & báo cáo giá trị nhóm \\
\bottomrule
\end{tabular}
\end{table}

% -----------------------------------------------------------------------------
\section{Tìm kiếm và báo cáo}
% -----------------------------------------------------------------------------

Các hàm báo cáo duyệt B-Tree qua \code{btTraverse} với một hàm gọi lại và một ``ngữ cảnh'' (\code{ctx}) chứa
điều kiện lọc và mảng kết quả. Khi cần trả về mảng có kích thước chưa biết trước (cảnh báo HSD, bán chạy), hàm
duyệt hai lượt: lượt đầu đếm, lượt sau cấp phát đúng kích thước rồi ghi.

\begin{itemize}
  \item \textbf{Tìm theo tên}: đổi tên và từ khóa sang chữ thường, giữ mặt hàng có tên chứa từ khóa, sắp theo
        tên. $O(n\cdot\ell + r\log r)$, $\ell$ là độ dài tên.
  \item \textbf{Liệt kê theo nhóm}: so sánh nhóm không phân biệt hoa thường (\code{"Do uong"} và
        \code{"DO UONG"} là một nhóm); để trống thì liệt kê tất cả, sắp theo nhóm rồi tên. $O(n\log n)$.
  \item \textbf{Cảnh báo tồn dưới ngưỡng}: \code{tonKho < nguongCanhBao}, sắp theo tồn tăng dần. $O(n + r\log r)$.
  \item \textbf{Cảnh báo HSD}: với mỗi lô còn hàng tính số ngày còn lại
        $d = \textsc{NgayTuyetDoi}(\mathit{hsd}) - \textsc{NgayTuyetDoi}(\mathit{homNay})$, giữ lô có $d \le 7$
        (kể cả $d < 0$: đã quá hạn nhưng vẫn còn hàng), sắp theo HSD. $O(n + L + r\log r)$.
  \item \textbf{Giá trị tồn theo nhóm}: không có \code{std::map}, nên gom tất cả mặt hàng, Merge Sort theo nhóm để các
        mặt hàng cùng nhóm nằm liền nhau, rồi quét một lượt cộng dồn số mặt hàng, tổng tồn và giá trị
        $\sum \mathit{soLuong}\times\mathit{giaNhap}$ của các lô. $O(n\log n + L)$.
\end{itemize}

\begin{algorithm}[htbp]
\caption{Top $k$ mặt hàng bán chạy trong tháng (\code{topBanChay})}\label{alg:top}
\begin{algorithmic}[1]
\Function{TopBanChay}{$\mathit{kho}, \mathit{thang}, \mathit{nam}, k$}
  \ForEach{mặt hàng $\mathit{mh}$ trong cây} $\mathit{mh}.\mathit{demTam} \gets 0$ \EndFor \Comment{$O(n)$}
  \ForEach{phiếu $p$ trong danh sách phiếu} \Comment{$O(P\log n)$}
    \If{$p.\mathit{loai} = \code{XUAT}$ \textbf{và} $p.\mathit{ngay}$ thuộc $\mathit{thang}/\mathit{nam}$}
      \State $\mathit{mh} \gets \Call{Tim}{\mathit{kho}.\mathit{cay}, p.\mathit{maVach}}$
      \If{$\mathit{mh} \ne \textsc{nil}$} $\mathit{mh}.\mathit{demTam} \gets \mathit{mh}.\mathit{demTam} + p.\mathit{soLuong}$ \EndIf \Comment{mã đã xóa: bỏ qua}
    \EndIf
  \EndFor
  \State $a \gets$ các mặt hàng có $\mathit{demTam} > 0$ (duyệt cây hai lượt)
  \State \Call{MergeSort}{$a$, giảm dần theo $\mathit{demTam}$, bằng nhau theo mã vạch}
  \State \Return $\min(k, |a|)$ phần tử đầu của $a$
\EndFunction
\end{algorithmic}
\end{algorithm}

Trường \code{demTam} được đặt lại về 0 ở đầu mỗi lần gọi nên gọi nhiều lần không bị cộng dồn.
Độ phức tạp $O(n + P\log n + r\log r)$.

% -----------------------------------------------------------------------------
\section{Xử lý ngày}
% -----------------------------------------------------------------------------

Ngày được quy đổi sang \emph{số ngày tuyệt đối} để cộng, trừ, so sánh trong $O(1)$. Coi tháng 1, 2 là tháng 13,
14 của năm trước (để ngày nhuận 29/02 nằm cuối ``năm''): với $y' = y - [m \le 2]$, $m' = m + 12[m \le 2]$,
\[
  N(d, m, y) = 365y' + \left\lfloor\tfrac{y'}{4}\right\rfloor - \left\lfloor\tfrac{y'}{100}\right\rfloor
  + \left\lfloor\tfrac{y'}{400}\right\rfloor + \left\lfloor\tfrac{153m' - 457}{5}\right\rfloor + d - 306 .
\]
Hàm ngược \code{daysToDate} khôi phục $(d, m, y)$ từ $N$. Khi đó
\begin{gather*}
  \mathit{congNgay}(D, k) = \mathit{daysToDate}\big(N(D) + k\big), \\
  \mathit{khoangCachNgay}(D_1, D_2) = N(D_1) - N(D_2).
\end{gather*}
Năm nhuận theo lịch Gregory: chia hết cho 4 nhưng không chia hết cho 100, hoặc chia hết cho 400. Ngày hợp lệ phải
có năm 1900--2100.

% -----------------------------------------------------------------------------
\section{Đọc và ghi tệp}
% -----------------------------------------------------------------------------

Khi khởi động, chương trình đọc lần lượt \code{mathang.txt}, \code{lo.txt}, \code{phieu.txt}. Mỗi dòng được bỏ
ký tự \code{\textbackslash r} cuối dòng (tệp tạo trên Windows), tách theo \code{|} và kiểm tra từng trường.
Dòng sai bị \emph{bỏ qua} và báo ra màn hình kèm tên tệp, số dòng và lý do, thay vì dừng chương trình.

\begin{algorithm}[htbp]
\caption{Đọc tệp lô (\code{docLo})}\label{alg:doclo}
\begin{algorithmic}[1]
\ForEach{dòng $s$ của \code{lo.txt}}
  \State tách $s$ thành 6 trường; \textbf{nếu} sai số trường: báo lỗi, bỏ qua
  \State $\mathit{mh} \gets \Call{Tim}{\mathit{kho}.\mathit{cay}, \mathit{maVach}}$ \Comment{mã vạch phải có trong \code{mathang.txt}}
  \State kiểm tra: mã lô dạng \code{LO}+số; hai ngày hợp lệ; $\mathit{hsd} \ge \mathit{ngayNhap}$; số lượng $> 0$; giá nhập $\ge 0$
  \State kiểm tra thứ tự FIFO: $\mathit{ngayNhap} \ge$ ngày nhập của lô cuối hàng đợi
  \State \Call{Enqueue}{$\mathit{mh}.\mathit{dsLo}$, lô mới}; $\mathit{mh}.\mathit{tonKho} \gets \mathit{mh}.\mathit{tonKho} + \mathit{soLuong}$ \Comment{tính lại tồn}
  \State $\mathit{kho}.\mathit{maLoKeTiep} \gets \max(\mathit{kho}.\mathit{maLoKeTiep}, \mathit{maLo} + 1)$
\EndFor
\end{algorithmic}
\end{algorithm}

Phiếu không bắt buộc mã vạch còn trong kho (phiếu của mặt hàng đã xóa vẫn là lịch sử), nhưng mã phiếu phải
tăng dần -- vừa phát hiện trùng mã trong $O(1)$ mỗi dòng, vừa khớp với cách ghi. Khi ghi, mỗi tệp được ghi ra
\code{<tệp>.tmp} trước rồi mới thay tệp cũ: nếu lỗi giữa chừng, dữ liệu cũ vẫn còn nguyên. Ở lần lưu đầu tiên
mỗi phiên, tệp cũ không bị xóa mà đổi tên thành \code{<tệp>.bak}, nên các dòng lỗi bị bỏ qua khi đọc (cùng các lô
của mặt hàng trên dòng đó) vẫn còn trong bản sao lưu để sửa tay. Đọc $O(n\log n + L\log n + P)$,
ghi $O(n + L + P)$.

% -----------------------------------------------------------------------------
\section{Tổng hợp độ phức tạp}
% -----------------------------------------------------------------------------

\begin{table}[htbp]
\centering\small
\caption{Độ phức tạp thời gian các thao tác chính}
\label{tab:dophuctap}
\begin{tabular}{@{}lll@{}}
\toprule
\textbf{Thao tác} & \textbf{Thời gian} & \textbf{Ghi chú} \\
\midrule
Tìm theo mã vạch & $O(\log_t n)$ nút, $O(\log n)$ so sánh & $h \le 1 + \log_t\frac{n+1}{2}$ \\
Thêm / xóa mặt hàng & $O(m\log_t n)$ & tách / mượn / gộp dọc một đường đi \\
Sửa mặt hàng & $O(\log n)$ & \\
Nhập kho & $O(\log n)$ & thêm lô cuối hàng đợi $O(1)$ \\
Xuất kho & $O(\log n + k)$ & $k$: số lô bị lấy hết \\
Tìm theo tên & $O(n\ell + r\log r)$ & duyệt toàn bộ cây \\
Liệt kê theo nhóm & $O(n\log n)$ & Merge Sort \\
Cảnh báo tồn / HSD & $O(n + L + r\log r)$ & \\
Giá trị tồn theo nhóm & $O(n\log n + L)$ & sắp theo nhóm rồi quét \\
Top 10 bán chạy & $O(n + P\log n + r\log r)$ & \code{demTam} + B-Tree \\
Đọc / ghi tệp & $O((n + L)\log n + P)$ / $O(n + L + P)$ & \\
\bottomrule
\end{tabular}
\end{table}
